%Faro, Tue Feb 11 23:19:56 WET 2014
%Cinematica e Din�mica do Ponto Material: 
%Movimento de um pendulo
%Written by TORDAR

\documentclass[12pt]{article}  

\pagestyle{empty}

\usepackage{graphicx}
\usepackage{pst-all} 
\usepackage{pst-slpe}

%Let's go: 

\begin{document}

\null\vskip4cm
\begin{center}
\begin{pspicture}(-0.1,-0.1)(1.1,1.1)
\psset{unit=8cm}
\psframe[linewidth=1pt,
         linecolor=lightgray,
	 dimen=inner,
	 fillstyle=hlines*,
	 fillcolor=lightgray,hatchangle=45](0.05,0.9)(0.95,1)
\psarc[linewidth=3pt,arcsepA=3pt,arcsepB=3pt]{->}(0.5,0.9){0.6}{270}{297}
\psline[linewidth=2pt,linecolor=black,linestyle=dashed,dash=5mm 1mm]{-}(0.5,0.9)(0.5,0)
\psline[linewidth=3pt,linecolor=black](0.5,0.9)(0.9,0.1)
\pscircle[linewidth=0.1mm,
linecolor=gray,
fillstyle=ccslope,
slopebegin=lightgray,
slopeend=black,
slopecenter=0.65 0.65,
slopesteps=25](0.9,0.1){0.05}
\psline[linewidth=3pt,linecolor=black](0,0.9)(1,0.9)
\rput{0}(1,0.07){\scalebox{1.5}{$m$}}
\rput{0}(0.65,0.25){\scalebox{1.5}{$\theta$}}
\end{pspicture}
\end{center}

\end{document} 
